\section{Introduction}

Suppose that we have many jobs waiting for one machine.  We want to minimize
the sum of the times at which they finish.  The length of each job is written
on a folded piece of paper.  If we knew all the lengths, the right schedule
would be clear: we would simply run the jobs from shortest to longest.  But
the papers are folded, and opening one of them takes one unit of machine time.

Should we open another paper, or should we start working?

Opening a paper may reveal a tiny job that we can finish immediately.  It may
also reveal a huge job that is best left for the end.  On the other hand,
looking has a cost.  If one hundred jobs are still unfinished, one unit spent
opening a paper adds one hundred units to the sum of completion times.  Thus
the machine has to balance two uses of its time: it can learn which jobs are
short, or it can complete jobs that are already understood.

This is the scheduling problem studied in this paper.  We start with its
cleanest version, in which every paper has to be opened.  We then allow a job
to be run without opening its paper.  Depending on what this means, we obtain
three natural variants.  The same simple ideas---thresholds, random sampling,
and shortest-first execution---give exact answers in all of them.

\subsection{First example: every job has to be tested}

There are $n$ jobs.  Testing job $i$ takes one unit of time and reveals a
number $p_i\geq0$.  The job then needs another $p_i$ units of processing.  A
job with $p_i=0$ is completed by its test.  For now, every job has to be
tested.

After a test, the scheduler faces one elementary decision.  Should the
revealed job be finished now, or should it be put aside while another job is
tested?  Short jobs should clearly be finished.  Long jobs should clearly
wait.  This suggests a threshold algorithm.

\begin{algorithm}[H]
\caption{The shape of an optimal obligatory-testing algorithm}
\label{alg:intro-obligatory}
\begin{algorithmic}
\State Test a job.
\State If its revealed length is below the current threshold, finish it now.
\State Otherwise put it aside and test another job.
\State After all tests, finish the jobs put aside from shortest to longest.
\end{algorithmic}
\end{algorithm}

The difficult part is choosing the threshold.  A deterministic algorithm has
to update it from the history it has seen.  A randomized algorithm can first
test a small random sample and use the sample to choose one threshold for the
rest of the input.  Some safeguards are needed when job lengths are
unbounded, but this is the whole algorithmic picture.

We compare the online scheduler with the shortest-first schedule that knows
all values $p_i$ in advance.  As usual, the competitive ratio is the online
cost divided by this offline optimum.  We are interested in the asymptotic
ratio: lower-order additive terms, of size $o(n^2)$, do not matter.

The obligatory-testing model for total completion time was introduced by
Dogeas, Erlebach, and Liang~\cite{DogeasErlebachLiang2024}.  Their work and a
later lower bound of Liang and Liang~\cite{LiangLiang2026} left a gap between
the best deterministic upper and lower bounds.  They also left the randomized
problem open.  We close both gaps.

\begin{theorem}[Obligatory testing, informal]
\label{thm:intro-obligatory}
The best deterministic asymptotic competitive ratio is
$1.5705740966\ldots$.  Against an adversary that fixes the input before the
algorithm draws its random bits, the best randomized ratio is $4/3$.

The randomized algorithm is not told the multiset of job lengths, and its
guarantee holds even when the lengths are unbounded.  This is the informal
version of \cref{thm:det-obligatory,thm:rand-obligatory}.
\end{theorem}

Let us briefly describe why the two constants are different.  A hard input
for a deterministic algorithm reveals positive jobs in decreasing levels and
hides many zero jobs until the end.  If the algorithm finishes the positive
jobs, it delays the zeros.  If it postpones them, it creates a costly tail.
The levels can be chosen so that these two losses balance at
$1.5705740966\ldots$.

For randomized algorithms the hard input is simpler.  Half of the jobs have
length zero and half have length two, and their labels are shuffled.  Until
the shuffle is uncovered, one unit of work completes only about half a job.
Counting the area under the number of unfinished jobs gives the lower bound
$4/3$.  A threshold learned from a random sample attains the same value.

\subsection{What if testing is optional?}

We now change one rule: a job may be run without a test.  There are two
separate questions to ask.

First, how long does an untested job run?  It may run for its true but hidden
length $p_i$.  Alternatively, every untested job may have the same known raw
length $u$, while a test replaces $u$ by a smaller value $p_i\leq u$.

Second, does the test reveal $p_i$?  If it does, the scheduler gains both a
faster job and information about where the job belongs.  If it does not, the
job becomes faster but remains unknown until it finishes.

These choices give the following four models.  The table is a summary rather
than four new sets of notation; in every row, the preliminary action in the
last column costs one unit of machine time.

\begin{center}
\small
\begin{tabular}{@{}L{0.20\linewidth}L{0.29\linewidth}L{0.40\linewidth}@{}}
\toprule
model & without a test & what the unit action does \\
\midrule
obligatory testing & the job cannot be run & reveals $p_i$ and gives access to the job \\
blind execution & the job runs for hidden time $p_i$ & reveals $p_i$ before we run it \\
common upper bound & the job runs for time $u$ & reveals and replaces $u$ by $p_i\leq u$ \\
blind optimization & the job runs for time $u$ & replaces $u$ by a still-hidden $p_i\leq u$ \\
\bottomrule
\end{tabular}
\end{center}

The last three rows separate the two benefits of the unit preliminary action.
Blind execution studies information without speedup.  Blind optimization
studies speedup without advance information.  The common-upper model contains
both.

\subsection{One algorithm for every bag}

Once testing becomes optional, a single worst-case ratio no longer tells the
whole story.  Consider the input as a bag $M$ of $n$ job lengths.  On one bag
it may be best to test almost every job.  On another bag it may be best to
test none.  We would like the algorithm to make the right choice on each bag,
not merely to survive the worst bag.

There is an obvious difficulty: the algorithm is not told $M$.  Still, it can
look at a small random sample.  Since the sample has size $o(n)$, it costs
only $o(n^2)$ in the objective.  Since it is random, it gives an accurate
picture of the whole bag.  The algorithm can then behave almost as if it had
known $M$ from the beginning.

This leads to a stronger goal than a competitive ratio.  For each bag $M$,
let $n^2\Phi(M)$ denote the best leading-order cost of an online policy
designed specifically for that bag.  We ask whether one policy, chosen before
seeing $M$, can attain $n^2\Phi(M)$ for every $M$.

\begin{theorem}[One algorithm for every bag, informal]
\label{thm:intro-instance-optimality}
Fix a bound on the job lengths.  In each of the three optional models there
is one randomized algorithm, independent of $M$, whose cost on every bag $M$
is
\[
                       n^2\Phi(M)+o(n^2).
\]
No online policy has a smaller leading cost on $M$, even if that competing
policy is told the whole bag in advance.  The three models have different,
explicit functions $\Phi$.

This is the informal version of
\cref{thm:common-upper-instance,thm:blind-instance,thm:bo-instance}.
\end{theorem}

The important point is the order of the quantifiers: the same algorithm
works for every bag.  Telling the bag to the competing policy is only a
convenient strengthening of the lower bound.  It is not extra information
available to our algorithm.

The randomization also has a simple role.  If jobs are presented under fixed
labels, a private random shuffle makes the pilot representative.  In the
equivalent random-bag picture, asking for a new job already returns a random
remaining element of the bag, and no additional randomization is needed
after the pilot.

The rest of the introduction explains what $\Phi(M)$ and its optimal policy
look like in the three models.  Bounded job lengths are needed for these
uniform instance-optimal statements.  They are not needed for the
randomized obligatory result above.

\subsection{Information without speedup: blind execution}

Suppose an untested job can simply be run until it finishes.  Its running
time is its true hidden length $p_i$.  Testing does not make the job shorter;
it only tells us $p_i$ before we decide when to run it.  We call this
\emph{blind execution}.

Why do we use the bag-dependent value $\Phi(M)$ here instead of the ordinary
clairvoyant optimum?  Take a bag containing many zero jobs and a few positive
jobs.  A scheduler that knows the bag's labeling completes all zeros at time
zero.  An online scheduler has to search for them.  Its competitive ratio can
therefore be arbitrarily large.  This does not mean that all online policies
are equally bad.  The quantity $\Phi(M)$ captures the best cost that online
information permits on this particular bag.

Let $\mu$ be the mean job length.  A revealed job shorter than $\mu$ should
be run before an untouched job; a longer one should be run after it.  When
testing is useful, there is another threshold $\tau$ that decides which jobs
are so short that they should interrupt the testing phase.  It is defined by
\[
                         \mathbb E(\tau-P)^+=1.
\]
The left-hand side is the expected saving obtained by moving outcomes below
$\tau$ forward.  Thus the cutoff is chosen exactly where this saving balances
the unit cost of a test.

One number remains: the fraction $q$ of jobs to test.  Once the bag is known,
the leading cost is an explicit quadratic in $q$, so the best value of $q$
is obtained by minimizing a one-variable quadratic.

\begin{algorithm}[H]
\caption{Instance-optimal blind execution}
\label{alg:intro-blind}
\begin{algorithmic}
\State Run a sublinear random pilot blindly and estimate the bag.
\State From the estimate compute $\mu$, $\tau$, and the tested fraction $q$.
\State Test a random $q$-fraction; finish every outcome below $\tau$ immediately.
\State Finish the remaining tested jobs below $\mu$, shortest first.
\State Run all untouched jobs blindly.
\State Finish the remaining tested jobs, shortest first.
\end{algorithmic}
\end{algorithm}

The block between testing and blind execution is important.  A job can be too
long to interrupt the test phase and still be shorter than an untouched job
on average.  Such a job waits until testing stops and is then processed before
the blind block.

A similar test-prefix structure is known in the Bayesian i.i.d. model of
Levi, Magnanti, and Shaposhnik
~\cite{LeviMagnantiShaposhnik2019,LeviMagnantiShaposhnik2024}.  Our input is
instead one adversarial bag.  The closest earlier adversarial model is the
processing-time oracle of Dufoss{\'e} et al.~\cite{DufosseEtAl2022}, where the
hidden length has only two possible values.  Our bag may contain arbitrary
bounded lengths.  The precise result is \cref{thm:blind-instance}.

\subsection{Information and speedup: a common upper bound}

We next assume that every untested job can be run for a known raw time $u$.
A unit test reveals a value $p_i\leq u$ and replaces the raw execution by the
shorter one.  Thus the test buys both information and speed.

The two endpoints are easy to understand.  If $u\leq1$, the test alone costs
at least as much as raw execution, so we should never test.  If $u$ is very
large, raw execution is unattractive, so the problem approaches obligatory
testing.  The interesting question is what happens between these endpoints.

This common-upper model was introduced by D{\"u}rr et
al.~\cite{DurrErlebachMegowMeissner2020}.  Their results already showed that
intermediate values of $u$ are the difficult ones.  We determine the exact
answer for every $u$, for both deterministic and randomized algorithms.

A deterministic algorithm uses three basic shapes.

\begin{algorithm}[H]
\caption{Deterministic common-upper policies, schematically}
\label{alg:intro-common-upper-det}
\begin{algorithmic}
\State For small $u$, run every job raw.
\State For intermediate $u$, test every job and use a fixed threshold.
\State For large $u$, use the history-dependent obligatory-testing threshold.
\State Whenever jobs are deferred, finish them shortest first.
\end{algorithmic}
\end{algorithm}

For randomized algorithms we can say more.  On each fixed bag, the same
sample-and-optimize principle from \cref{thm:intro-instance-optimality}
finds the best threshold and the best tested fraction $q$.

\begin{algorithm}[H]
\caption{Instance-optimal common-upper policy}
\label{alg:intro-common-upper}
\begin{algorithmic}
\State Test a sublinear random pilot and estimate the bag.
\State Compute the best threshold and the best tested fraction $q$.
\State Test a random $q$-fraction; finish outcomes below the threshold at once.
\State Finish all other tested jobs, shortest first.
\State Run every untouched job raw for time $u$.
\end{algorithmic}
\end{algorithm}

For each bag, this policy attains its best online leading cost.  We obtain the
competitive ratio by asking which bag maximizes the ratio between that cost
and the clairvoyant optimum.  The worst bag can be reduced to two job
lengths, leaving a one-variable optimization.

\begin{theorem}[Common-upper curves, informal]
\label{thm:intro-common-upper-curves}
For every $u>0$, the exact deterministic and randomized asymptotic
competitive ratios are the six-piece and four-piece curves in
\cref{fig:intro-curves}.  The deterministic curve reaches
$1.866760\ldots$ and then falls to the obligatory value
$1.570574\ldots$.  The randomized curve reaches $1.625752\ldots$ and
eventually becomes $4/3$.

This is the informal version of
\cref{thm:optional-curve,thm:random-common-upper}.
\end{theorem}

\begin{figure}[H]
\centering
% The introduction curves as native pgfplots graphics.  All explicit
% branches are plotted from their formulas.  The only tabulated branch is
% C_mix, which is defined implicitly in Equation (intro-mixed-system).
\begingroup
% TeX Live 2023's pgfplots emits spurious nullfont diagnostics containing
% ``0.01pt'' while sampling expressions; no text is actually lost.
\tracinglostchars=0
\pgfmathdeclarefunction{introRhoI}{1}{%
  \pgfmathparse{%
    (-(((#1-1)^3)+((#1-1)^2)+1)
      +sqrt((((#1-1)^3)+((#1-1)^2)+1)^2
        +4*(#1-1)*(((#1-1)^2)+(#1-1)+1)*((#1-1)+2)))
      /(2*(#1-1))}%
}
\begin{tikzpicture}
\begin{groupplot}[
  group style={group size=2 by 1,horizontal sep=1.15cm},
  width=0.465\textwidth,
  height=0.315\textwidth,
  xmin=0.2,
  xmax=7.3,
  ymin=0.97,
  ymax=1.91,
  xlabel={raw execution time $u$},
  xtick={1,2,3,4,5,6,7},
  ytick={1,1.2,1.4,1.6,1.8},
  tick label style={font=\scriptsize},
  label style={font=\small},
  title style={font=\small},
  grid=major,
  grid style={black!15},
  axis line style={black!65},
  tick style={black!65},
  axis x line*=bottom,
  axis y line*=left,
  samples=90,
  no markers,
  clip=false,
]

\nextgroupplot[
  title={deterministic},
  ylabel={optimal competitive ratio},
]
\addplot[blue!72!black,very thick,domain=0.2:1] {1};
\addplot[blue!72!black,very thick,domain=1:1.86676039917] {x};
\addplot[blue!72!black,very thick,domain=1.86676039917:3.14789903570]
  {introRhoI(x)};
\addplot[blue!72!black,very thick,domain=3.14789903570:3.61803398875]
  {1+1/sqrt(x-1)};
\addplot[blue!72!black,very thick,smooth] coordinates {
  (3.61803399,1.61803399)
  (3.63837920,1.61566099)
  (3.65923107,1.61328800)
  (3.68063111,1.61091500)
  (3.70262731,1.60854201)
  (3.72527568,1.60616902)
  (3.74864222,1.60379602)
  (3.77280560,1.60142303)
  (3.79786084,1.59905003)
  (3.82392445,1.59667704)
  (3.85114178,1.59430404)
  (3.87969805,1.59193105)
  (3.90983519,1.58955805)
  (3.94187917,1.58718506)
  (3.97628649,1.58481206)
  (4.01372962,1.58243907)
  (4.05526944,1.58006608)
  (4.10275400,1.57769308)
  (4.15995362,1.57532009)
  (4.23729031,1.57294709)
  (4.50524150,1.57057410)
};
\addplot[blue!72!black,very thick,domain=4.50524149579:7.3]
  {1.57057409665};
\addplot[black!45,densely dotted,thin] coordinates {(1,0.97) (1,1.91)};
\addplot[black!45,densely dotted,thin]
  coordinates {(1.86676039917,0.97) (1.86676039917,1.91)};
\addplot[black!45,densely dotted,thin]
  coordinates {(3.14789903570,0.97) (3.14789903570,1.91)};
\addplot[black!45,densely dotted,thin]
  coordinates {(3.61803398875,0.97) (3.61803398875,1.91)};
\addplot[black!45,densely dotted,thin]
  coordinates {(4.50524149579,0.97) (4.50524149579,1.91)};
\addplot[only marks,mark=*,mark size=1.45pt,black]
  coordinates {(1.86676039917,1.86676039917) (4.50524149579,1.57057409665)};
\node[font=\scriptsize,anchor=west] at (axis cs:2.25,1.845)
  {maximum $1.87\ldots$};
\draw[black!55,thin] (axis cs:2.21,1.85) -- (axis cs:1.89,1.86676);
\node[font=\scriptsize,anchor=west] at (axis cs:5.05,1.605)
  {$1.57\ldots$};

\nextgroupplot[title={randomized}]
\addplot[orange!88!black,very thick,domain=0.2:1] {1};
\addplot[orange!88!black,very thick,domain=1:5.04891733952]
  {x^3/(x^2+(x-1)^3)};
\addplot[orange!88!black,very thick,domain=5.04891733952:6.25]
  {1+1/(2*(sqrt(x)-1))};
\addplot[orange!88!black,very thick,domain=6.25:7.3] {4/3};
\addplot[black!45,densely dotted,thin] coordinates {(1,0.97) (1,1.91)};
\addplot[black!45,densely dotted,thin]
  coordinates {(5.04891733952,0.97) (5.04891733952,1.91)};
\addplot[black!45,densely dotted,thin]
  coordinates {(6.25,0.97) (6.25,1.91)};
\addplot[only marks,mark=*,mark size=1.45pt,black] coordinates {
  (2.36602540378,1.62575238458)
  (5.04891733952,1.39719767662)
  (6.25,1.33333333333)
};
\node[font=\scriptsize,anchor=west] at (axis cs:2.88,1.65)
  {maximum $1.63\ldots$};
\draw[black!55,thin] (axis cs:2.84,1.647) -- (axis cs:2.39,1.627);
\node[font=\scriptsize,anchor=west] at (axis cs:6.53,1.375) {$4/3$};

\end{groupplot}
\end{tikzpicture}
\endgroup

\caption{The exact asymptotic ratios as a function of the raw execution time
$u$.  Testing is pointless on the left.  On the right, both curves reach
their obligatory-testing values.}
\label{fig:intro-curves}
\end{figure}

The curves are not monotone.  This is not a contradiction: increasing $u$
changes both players' options.  For intermediate $u$, the clairvoyant
scheduler makes especially good use of the choice between raw execution and
testing.  When $u$ becomes large, both schedulers effectively test
everything.

The hard bags change along the curves.  In the first regimes, mixtures of
zero jobs and jobs of length $u$ are enough.  Later, the deterministic lower
bound becomes a capped form of the decreasing-level construction from
obligatory testing.  The exact formulas and matching constructions are in
\cref{sec:optional,sec:random-common-upper}.

\subsection{Speedup without information: blind optimization}

Finally, suppose that a unit-time optimization replaces $u$ by $p_i\leq u$
but does not reveal $p_i$.  The value becomes known only when the optimized
job finishes.  Now a threshold cannot separate short and long outcomes:
before execution, all optimized jobs still look the same.

This makes the optimal policy simpler.  A pilot estimates only the mean
$\mu$ of the optimized lengths.  Afterwards, the scheduler compares the two
possible expected lengths, $u$ and $1+\mu$, and uses the better action for all
remaining jobs.

\begin{algorithm}[H]
\caption{Instance-optimal blind optimization}
\label{alg:intro-blind-optimization}
\begin{algorithmic}
\State Optimize and run a sublinear random pilot; estimate its mean $\widehat\mu$.
\State If $1+\widehat\mu<u$, optimize and run every remaining job.
\State Otherwise run every remaining job raw for time $u$.
\end{algorithmic}
\end{algorithm}

This policy is again instance optimal up to $o(n^2)$.  Its worst-case ratio,
however, grows as $\sqrt u/2$.  In contrast, the revealing common-upper curve
eventually stays at $4/3$.  The difference shows exactly what information is
worth: speedup alone does not give a bounded ratio when $u$ grows.

\begin{theorem}[Blind-optimization curve, informal]
\label{thm:intro-bo-curve}
The exact randomized asymptotic competitive ratio is one for $u\le1$ and is
given by an explicit two-branch curve for $u>1$.  After the transition point
$u=2.246979\ldots$, it equals
\[
 \frac12\left(1+\sqrt{\frac{u^2+u-1}{u-1}}\right).
\]
In particular, it is $\sqrt u/2+1/2+o(1)$ as $u\to\infty$, and hence is
unbounded.  This is the informal version of
\cref{thm:bo-curve}.
\end{theorem}

A related non-revealing optimization appears in the oblivious-budget model
of Damerius et al.~\cite{DameriusEtAl2023}.  There, optimization consumes an
external budget.  Here it occupies the machine and delays every unfinished
job.  Our exact result is in \cref{sec:blind-optimization}.

\subsection{Techniques}

The proofs repeatedly use three pictures.

\paragraph{The area of unfinished jobs.}
During one unit of machine work, the objective increases by the number of
jobs that are still unfinished.  Thus the sum of completion times is the area
under the curve ``number of unfinished jobs versus time.''  This picture is
often easier to reason about than individual completion times.

\paragraph{Jobs completed per unit of work.}
A round of tests together with the short jobs completed during that round can
be viewed as one large action.  It spends some amount of time and completes
some fraction of the jobs.  Higher-density actions---those completing more
jobs per unit of work---should come first.  Sorting these actions produces
the thresholds, the ordered blocks, and the one-variable quadratics in the
instance-optimal results.

\paragraph{A private shuffle.}
The universal algorithms start by shuffling the labels and inspecting a
small prefix.  The same idea is needed in the lower bound, where the policy
may adaptively decide which label to touch next.  If equal job lengths are
temporarily distinguished as separate tokens, the order in which any
deterministic policy first reveals these tokens is itself just a permutation
of them.  The revealed sequence is therefore sampling without replacement.
One concentration event can then control every possible stopping point of
the policy.

The deterministic proofs use a complementary idea.  Their lower bounds make
the algorithm choose between completing the current jobs and waiting for
shorter future jobs.  Their upper bounds maintain a threshold and charge each
wrongly ordered pair to a small potential.  The exact constants appear when
the two possible losses are equal.

\subsection{Related work}
\label{sec:related}

The closest papers have already appeared next to the models they introduced.
We give a broader comparison here.

Scheduling with tests belongs to optimization with \emph{explorable
uncertainty}, where uncertain data can be queried at a cost; the query
viewpoint goes back at least to Kahan~\cite{Kahan1991}.  Levi, Magnanti, and
Shaposhnik study a stochastic scheduling model and derive prefix-testing
policies followed by an order based on expected processing time per unit
weight~\cite{LeviMagnantiShaposhnik2019,LeviMagnantiShaposhnik2024}.  Our
four-block blind-execution policy has a similar shape, but our input is an
adversarial bag, need not converge to a distribution, and is not announced to
the universal algorithm.

D{\"u}rr et al.~\cite{DurrErlebachMegowMeissner2020} introduced the
adversarial common-upper model and obtained deterministic bounds, including
the $1.8668\ldots$ guarantee for binary values.  We close the deterministic
gap for arbitrary values $p_i\in[0,u]$ and give the exact randomized curve
against an oblivious adversary.  Related work allows nonuniform test
times~\cite{AlbersEckl2021,LiuLiuWongZhang2023,KrekelbergEtAl2026}, multiple
machines~\cite{GongChenHayashi2024}, or weights~\cite{BuldSchulz2026}.

Dogeas, Erlebach, and Liang~\cite{DogeasErlebachLiang2024} introduced
obligatory testing for total completion time.  For uniform tests they proved
the SIDLE upper bound $1.58451\ldots$ and the deterministic lower bound
$\sqrt2$.  Liang and Liang~\cite{LiangLiang2026} later raised the lower bound
to $3/2$.  Our matching value is $1.5705740966\ldots$, and our randomized
result gives the oblivious value $4/3$.

The processing-time oracle of Dufoss{\'e} et al.~\cite{DufosseEtAl2022} is
close to blind execution, but allows only two possible hidden durations.  If
no tests are available and processing times are revealed only at completion,
one obtains the classical nonclairvoyant model of Motwani, Phillips, and
Torng~\cite{MotwaniPhillipsTorng1994}.  These semantics differ from raw
common-upper execution, which always lasts the declared time $u$.

The rest of the paper gives precise versions and proofs.  We first fix the
models and offline benchmarks in \cref{sec:model}.  Each main section then
starts with the theorem it proves and a short overview of the proof.  The
machine-checked part of the argument is summarized in
\cref{sec:formalization}.
